\documentclass{misisarticle}

\grnti{}
\asjc{}
\oecd{1.02}

\title{Классификация болезней и стрессов растений методами глубокого обучения}

\misisauthor{Осиновская Е. Д.}{Кафедра инженерной кибернетики, НИТУ «МИСиС»}{Москва, Россия}{student@edu.misis.ru}

\setabstract{Рассматривается задача многоклассовой классификации болезней растений по изображениям. Улучшается точность и устойчивость классификации при дисбалансе классов и ограниченном объеме данных. Применяются методы глубокого обучения (ResNet50, DenseNet121 с SE-блоками, InceptionV3, ViT-Base/Small), трансферное обучение, прогрессивная разморозка, аугментации (MixUp, TTA), функции потерь с весами классов (Focal Loss, Label Smoothing) и мета-классификатор для ансамблирования. Проверка проводится на двух открытых датасетах с оценкой точности, F1-score и визуализацией эмбеддингов (UMAP).}
\setkeywords{глубокое обучение, классификация изображений растений, трансферное обучение, несбалансированные данные, ансамбли моделей, болезни растений}

\begin{document}
\maketitle

\section{Введение}
Современное сельское хозяйство сталкивается с рядом вызовов, связанных с необходимостью повышения урожайности и обеспечения продовольственной безопасности. Одним из ключевых факторов, влияющих на продуктивность, являются заболевания и стрессовые состояния растений. Своевременное обнаружение этих проблем критически важно для минимизации потерь. В последние годы глубокое обучение продемонстрировало значительный потенциал в задачах классификации изображений, в том числе в агрономической сфере. Данная работа посвящена исследованию применения методов глубокого обучения для классификации болезней и стрессов растений по изображениям листьев. 

Вклад работы состоит в сравнительном исследовании свёрточных (ResNet50, DenseNet121 с SE-блоками, InceptionV3) и трансформерных (ViT-Base/Small) архитектур для задачи классификации болезней и стрессов растений на двух разнородных датасетах с разным уровнем дисбаланса классов. В работе единообразно оценивается влияние методов регуляризации, аугментации и функций потерь (Class Weights, Focal Loss, Label Smoothing, MixUp, TTA) на устойчивость моделей к дисбалансу и ограниченному объёму данных. Дополнительный вклад — построение мета-классификатора на основе логистической регрессии, объединяющего предсказания DenseNet121 + SE, InceptionV3 и ViT-Small.

\section{Анализ работ}
% Здесь будет текст анализа работ (пока оставлено пустым, как вы просили)
% Можно использовать следующие источники для вставки в этот раздел:
% 1. Ососков, Г. А. Применение методов глубокого обучения для решения различных задач в сельском хозяйстве / Г. А. Ососков, А. В. Ужинский, А. В. Нечаевский. // Государственный университет "Дубна". 30 лет в науке : сборник научных трудов. — Дубна : Университет Дубна, 2024. — С. 289-302.
% 2. Кутырёв, А. И. Распознавание и классификация болезней листьев яблони на основе анализа их изображений моделями сверточных нейронных сетей (CNN) / А. И. Кутырёв. // Вестник Ульяновской государственной сельскохозяйственной академии. — 2023. — № 3 (63). — С. 215-223.
% 3. Кондратьев, А. А. Применение нейронных сетей для обнаружения аномалий в развитии растений / А. А. Кондратьев, С. Б. Березин. // International Journal of Open Information Technologies. — 2023. — № 7. — С. 47-52.
% 4. Елисеев, А. Н. Решение задачи классификации сельскохозяйственных культур с использованием глубоких нейронных сетей / А. Н. Елисеев, И. И. Курочкин. // Облачные и распределенные вычислительные системы в электронном управлении. ОРВС-2023 : сборник трудов. —, 2024. — С. 35-40.
% 5. Демин, Н. И. Нейросетевые методы для распознавания фруктов и овощей / Н. И. Демин, В. С. Матвеев. // Искусственный интеллект в промышленных, коммерческих, медицинских и финансовых приложениях: сборник статей научно-технологического семинара кафедры "Инженерной кибернетики". — М. : Национальный исследовательский технологический университет "МИСИС", 2026. — С. 34-39.
% 6. Егоров, Ф. Д. Распознавание дефектов картофеля с помощью свёрточной нейронной сети / Ф. Д. Егоров. // Искусственный интеллект в промышленных, коммерческих, медицинских и финансовых приложениях: сборник статей научно-технологического семинара кафедры "Инженерной кибернетики". — М. : Национальный исследовательский технологический ун-т "МИСИС", 2026. — С. 40-45.

\section{Постановка задачи}
Целью данной научно-исследовательской работы является сравнительный анализ эффективности различных архитектур нейронных сетей и методов обучения в задаче многоклассовой классификации заболеваний и стрессовых состояний растений. Для достижения этой цели были поставлены следующие задачи:
\begin{enumerate}
    \item Провести разведочный анализ двух открытых датасетов: заболеваний рисовых листьев и стрессовых состояний растений.
    \item Оценить эффективность свёрточных нейронных сетей (ResNet50, DenseNet121 с SE-блоками, InceptionV3) и трансформерных архитектур (ViT-Base, ViT-Small) с использованием стратегий трансферного обучения и прогрессивной разморозки слоёв.
    \item Исследовать влияние методов регуляризации и аугментации данных, таких как MixUp, TTA (Test Time Augmentation) и Focal Loss, на качество классификации.
    \item Разработать мета-классификатор для ансамблирования предсказаний лучших моделей с целью повышения итоговой точности.
\end{enumerate}

\section{Описание методов}
В работе используются два открытых датасета. Первый датасет содержит изображения рисовых листьев, разделённые на 7 классов (blast, healthy, insect, leaf\_folder, scald, stripes, tungro). Второй датасет включает 4749 изображений высокого разрешения, сгруппированных в 10 классов, охватывающих 8 сельскохозяйственных культур, с признаками грибковых заболеваний, повреждений насекомыми и дефицита питательных веществ.

Для классификации применяются следующие архитектуры:
\begin{itemize}
    \item \textbf{ResNet50}: глубокая свёрточная сеть с остаточными соединениями (skip-connections), позволяющая эффективно обучать модели большой глубины.
    \item \textbf{DenseNet121 с SE-блоками}: свёрточная сеть с плотными соединениями между слоями для эффективной передачи градиентов. Блоки Squeeze-and-Excitation (SE) адаптивно перекалибровывают каналы признаков, выделяя наиболее важные.
    \item \textbf{InceptionV3}: свёрточная сеть, использующая Inception-модули с параллельными свёртками разных масштабов и факторизацию свёрток для захвата как мелких деталей, так и крупных паттернов.
    \item \textbf{Vision Transformer (ViT-Base и ViT-Small)}: трансформерная архитектура, рассматривающая изображение как последовательность патчей и использующая механизм Self-Attention для учёта глобальных взаимосвязей.
\end{itemize}

В качестве стратегии обучения используется \textbf{Transfer Learning} с предобучением на ImageNet. Для адаптации моделей применяется стратегия \textbf{Progressive Unfreezing} — постепенная разморозка слоёв от классификатора к более глубоким слоям. В экспериментах используются оптимизаторы Adam, AdamW и RAdam, планировщики StepLR и Linear warmup с CosineAnnealing. Для борьбы с дисбалансом классов применяются функции потерь CrossEntropyLoss с весами классов, Focal Loss и LabelSmoothingCrossEntropy. Для повышения обобщающей способности используются Gradient Clipping, MixUp и TTA. Ансамблирование предсказаний трёх моделей (DenseNet121+SE, InceptionV3, ViT-Small) выполняется с помощью мета-классификатора на основе логистической регрессии.

\section*{Литература}
% Здесь будет список литературы (пока оставлено пустым, как вы просили)
% В шаблоне список литературы оформляется через окружение thebibliography:
% \begin{thebibliography}{99}
% \bibitem{label1} Ососков, Г. А. ...
% \end{thebibliography}

\end{document}
